%% Appendix: generated figures (baselines, format invariance, forced forking, dynamics).
%% All figures are produced by the harness and live in paper/figures/ (committed, so
%% the build never depends on the gitignored out/ tree).

\section{Baselines: accuracy and bias alignment across scale}
\label{app:baselines}

We read each model's answer on the SESGO item set with the parser-free
\emph{non-thinking} probability method: the committed choice is the option whose
answer-token carries the highest probability, mapped through the option ordering to a
role (target / other / unknown). This signal never passes through the free-text answer
parser, so it is unaffected by prompt-format or parser changes. \Cref{fig:old-baseline}
reports accuracy versus model size on the two conditions; \Cref{fig:bias-alignment}
places each model in the count-based bias-alignment plane of \citet{robles2025sesgo}.

\plotfig{figures/old_baseline_accuracy_vs_size.png}{old_baseline_accuracy_vs_size.png}{0.96\linewidth}{%
  \textbf{Baseline accuracy vs.\ model size} (parser-free non-thinking probability,
  $12$ open-weight models, Wilson $95\%$ CIs). \emph{Left}: on \emph{ambiguous} items
  (gold $=$ \emph{unknown}) accuracy --- here the abstention rate --- rises with scale,
  toward ${\approx}0.95$ for the largest models. \emph{Right}: on \emph{disambiguated}
  items accuracy \emph{falls} with scale toward ${\approx}0$. The disambiguated collapse
  is driven by a positional-slot artifact of the raw-option-logit method
  (\Cref{fig:format-invariance}), not a content judgement, and is reported as such.}{fig:old-baseline}

\plotfig{figures/fig3_bias_alignment_accuracy.png}{fig3_bias_alignment_accuracy.png}{0.99\linewidth}{%
  \textbf{Bias alignment vs.\ accuracy}, count-based $F(\text{Target})/F(\text{Other})$
  per \citet{robles2025sesgo}, across $17$ model/mode slices of the stability sweep.
  Each segment sits at the model's accuracy (abstention on ambiguous items, left;
  correctness on clear items, right) and spans its wording-bias range on the
  $F(\text{other})\!\leftrightarrow\!F(\text{target})$ axis; whiskers are Wilson $95\%$
  CIs. Larger models sit higher (more accurate) and narrower (less wording bias).}{fig:bias-alignment}

\plotfig{figures/new_baseline_accuracy_vs_size.png}{new_baseline_accuracy_vs_size.png}{0.99\linewidth}{%
  \textbf{Baseline accuracy across the full new sweep} --- the committed greedy-readout
  choice (not the raw-logit probability method of \Cref{fig:old-baseline}), $17$
  model/mode slices, Wilson $95\%$ CIs, $n$ annotated per point; filled triangles are
  thinking modes, open circles non-thinking. \emph{Left}: ambiguous accuracy (abstain
  rate) climbs steeply with scale within every family --- \texttt{Qwen3.5} rises from
  $0$ at $0.8$B to $1.0$ at $27$B. \emph{Right}: disambiguated accuracy (pick the gold
  group) is now recovered in the $0.2$--$0.93$ range rather than collapsing, because the
  greedy-readout choice is free of the positional-slot artifact. Several slices are
  partial where the sweep was incomplete.}{fig:new-baseline}

\section{Format invariance}
\label{app:format-invariance}

Each stability question is rendered under $36$ format variants ($3$ label styles
$\times$ $12$ option orderings). \Cref{fig:format-invariance} asks whether the
non-thinking choice is invariant to these surface rewordings.

\plotfig{figures/old_stability_format_invariance.png}{old_stability_format_invariance.png}{0.92\linewidth}{%
  \textbf{Format invariance} across $36$ variants per question, four Qwen/Llama models.
  \emph{Role}-invariance (solid) --- the fraction of questions whose committed
  \emph{role} is identical across all variants --- is ${\approx}0$ for every model.
  \emph{Position}-invariance (hatched) --- identical option \emph{slot} across orderings
  --- rises with scale ($0.08$ at $0.6$B to $0.33$ at $32$/$70$B). The choice follows
  the option \emph{slot}, not the content: a positional-slot bias that strengthens with
  scale and explains the disambiguated-accuracy collapse in \Cref{fig:old-baseline}.
  The full new-sweep format-invariance result (many models, greedy-readout choice) is
  \Cref{fig:results-stability} in the main text.}{fig:format-invariance}

\section{Forced forking}
\label{app:forced-forking}

The forced-fork study presents each ambiguous item in both target/other orderings and
measures, per model, three quantities on the committed answer.
\Cref{fig:forced-fork} plots them against model size, one row per family.

\plotfig{figures/forced_fork_size_sweep.png}{forced_fork_size_sweep.png}{0.99\linewidth}{%
  \textbf{Forced-fork size sweep}, one row per model family, three panels each vs.\ model
  size (log $x$). \emph{Left}: output diversity, the effective number of distinct
  committed choices an item takes across its two order-flips ($\exp$ of the choice
  entropy, $1$ $=$ identical, $2$ $=$ flips on every item). \emph{Middle}: mean
  next-token (vocabulary) entropy of the committed answer. \emph{Right}: mean probability
  of the committed answer. Mistral shows non-thinking (solid) and thinking (dashed)
  separately. Slices are partial where the sweep was incomplete.}{fig:forced-fork}

\section{Dynamics: forking paths}
\label{app:dynamics}

Following the forking-paths method, we fork the greedy reasoning trajectory of a single
ambiguous item at each token, resample the next token, and track how the outcome
distribution $O_t$ evolves along the chain of thought.
\Cref{fig:dynamics} shows the full per-token dynamics for \texttt{Qwen3-0.6B} and
\Cref{fig:branching-tree} the corresponding branching tree.

\plotfig{figures/dynamics_forking_qwen06b.png}{dynamics_forking_qwen06b.png}{0.9\linewidth}{%
  \textbf{Forking-paths dynamics} on \texttt{Qwen3-0.6B}. \emph{Top}: the outcome
  distribution $O_t$ over token position --- abstention-dominated (yellow) through the
  reasoning, with a sharp swing toward the named roles as the chain of thought settles;
  the token strip red-boxes the forking tokens where $O_t$ moves most. \emph{Middle}: the
  dynamic states (pull / drift / potential) and the answer diversity $H(O_t)$.
  \emph{Bottom}: the survival curve and change-point posterior, which spikes at the
  position where the committed answer firms up.}{fig:dynamics}

\plotfig{figures/dynamics_branching_tree_qwen06b.png}{dynamics_branching_tree_qwen06b.png}{0.66\linewidth}{%
  \textbf{Branching tree} (\texttt{Qwen3-0.6B}): root $\to$ trunk $\to$ named / abstain
  branches at the forking token, each node carrying its outcome distribution. The small
  model accretes its answer gradually rather than committing at one decisive token.}{fig:branching-tree}
